\chapter{Characterisation of Clusters}\label{app:candidate_characterisation}

\section{Bayesian Model Details}\label{sec:app_ch6_bayesian_details}

Section~\ref{sec:ch6_bayesian_characterisation} defines the model structure, predictor scales, analytical levels, and coefficient interpretation. The Gaussian models approximate the bounded 0--1 screening scores and use a common residual standard deviation within each fit. This appendix provides the fitted sample sizes, priors, sampling settings, and diagnostics.

All 8,962 size-eligible clusters entered the CSC and local burst models, and all 4,216 burden-eligible clusters entered the onward burden models. Expanded composition models excluded five additional sequence records for CSC and local burst outcomes and three for onward burden score because of missing surveillance measures. Table~\ref{tab:bayesian_model_specifications} gives the sample sizes for all 18 models.

\begingroup
    \let\thesistablesetup\thesisdensetablesetup
    \input{assets/scotland/tab_bayesian_model_specifications}
\endgroup

Models were implemented in Bambi using PyMC \citep{Capretto2022,AbrilPla2023}, with the explicitly specified priors in Table~\ref{tab:bayes_priors}. Intercept priors are centred on the fitted sample's outcome mean (or its logit for CSC). Bambi centres the fixed-effect design columns internally during sampling, so this prior applies to the intercept in that centred parameterisation; reported intercepts are transformed back to the original predictor coding. We used non-centred parameterisation for varying intercepts and four sampling chains, each with 2,000 tuning iterations followed by 2,000 retained posterior draws (8,000 draws in total), target acceptance 0.99, and random seed 123.

\begin{table}[htbp]
    \centering
    \caption[Priors for Bayesian characterisation models]{Priors specified for Bayesian characterisation models. Here, $\bar y$ is the observed outcome mean in the data used to fit each model. The intercept prior refers to the internally centred fixed-effect design. Normal distributions are parameterised by mean and standard deviation; half-normal distributions specify priors on standard deviations. Separate group-level standard deviations are estimated for policy period and clade.}\label{tab:bayes_priors}
    \begin{thesistablebody}{@{}P{0.26\textwidth}P{0.35\textwidth}P{0.35\textwidth}@{}}
    \toprule
    \textbf{Term} & \textbf{Logistic model} & \textbf{Linear model} \\
    \midrule
    Intercept & $\operatorname{Normal}\{\operatorname{logit}(\bar y),1.5\}$ & $\operatorname{Normal}(\bar y,1.0)$ \\
    Fixed effects & $\operatorname{Normal}(0,1.0)$ & $\operatorname{Normal}(0,0.5)$ \\
    Varying intercepts & $\operatorname{Normal}\{0,\operatorname{HalfNormal}(1.0)\}$ & $\operatorname{Normal}\{0,\operatorname{HalfNormal}(0.5)\}$ \\
    Residual SD & Not applicable & $\operatorname{HalfNormal}(0.5)$ \\
    \bottomrule
    \end{thesistablebody}
\end{table}

Odds-ratio means and HDIs are calculated from exponentiated posterior draws. Exponentiated logistic intercepts represent baseline odds, while exponentiated period and clade deviations are odds multipliers relative to the baseline with no group deviation. Exponentiated group-level standard deviations describe multiplicative variation on the odds scale, rather than standard deviations of odds ratios. Table~\ref{tab:app_bayes_diagnostic_guide} defines the sampler diagnostics and warning thresholds used in Table~\ref{tab:bayesian_model_diagnostics}. These assess sampling reliability; checks of how well the models reproduce observed data (posterior predictive checks) are not reported.

\begin{table}[htbp]
    \centering
    \caption[Guide to Bayesian sampler diagnostic metrics]{Guide to Bayesian sampler diagnostic metrics and warning thresholds for the fitted characterisation models. Definitions follow diagnostic guidance for the No-U-Turn Sampler (NUTS)\nomenclature{NUTS}{No-U-Turn Sampler} and Hamiltonian Monte Carlo (HMC)\nomenclature{HMC}{Hamiltonian Monte Carlo} and rank-normalised $\widehat R$/ESS recommendations \citep{HoffmanGelman2014,Betancourt2017,Vehtari2021,StanDevelopmentTeam2025}. BFMI denotes the Bayesian fraction of missing information\nomenclature{BFMI}{Bayesian fraction of missing information}.}\label{tab:app_bayes_diagnostic_guide}
    \begin{thesistablebody}{@{}P{0.16\textwidth}P{0.38\textwidth}P{0.40\textwidth}@{}}
    \toprule
    \textbf{Metric} & \textbf{What it checks} & \textbf{Interpretation used here} \\
    \midrule
    Divergences & Post-warmup NUTS transitions with unreliable Hamiltonian integration. & Required: 0. Any divergence is a validity warning. \\
    Min BFMI & Lowest chain-level BFMI; checks marginal energy exploration after momentum resampling. & Required: $\geq0.3$. Lower values indicate poor energy exploration. \\  %chktex 13
    Max $\widehat R$ & Largest rank-normalised split-$\widehat R$; compares within- and between-chain variation. & Required: $\leq1.01$. Larger values suggest incomplete mixing or non-stationarity. \\  %chktex 13
    Min bulk ESS & Smallest bulk ESS; effective number of independent draws for the central part of the posterior distribution. & Required: $\geq400$. Lower values weaken mean and median summaries. \\  %chktex 13
    Min tail ESS & Smallest tail ESS; effective draws for posterior tails and interval endpoints. & Required: $\geq400$. Lower values weaken HDIs and tail probabilities. \\  %chktex 13
    Max tree depth & Largest NUTS binary-tree depth reached when choosing a trajectory length. & Recorded for context; no automatic warning threshold is applied in the diagnostic classification. \\ %chktex 13
    \bottomrule
    \end{thesistablebody}
\end{table}

\input{assets/scotland/tab_bayesian_model_diagnostics}

\clearpage
\section{Cluster Size, Screening Route, and Descriptive Sensitivity}\label{sec:app_ch6_size_sensitivity}

These summaries use saved screening outputs and the same retained windows as the primary analysis. They do not involve Bayesian refitting or screening recalibration. The primary denominator is 8,962 clusters and 152,177 sequence records, with complete values for the five composition attributes and the six displayed entropy measures. Table~\ref{tab:ch6_size_overview} in the main chapter gives denominators by size band; Table~\ref{tab:app_ch6_size_composition} below gives the corresponding composition percentages and contrasts among pooled records.

\begin{figure}[htbp]
    \centering
    \includegraphics[width=\textwidth]{fig_app_ch6_composition_by_size}
    \caption[Composition across size bands with equal cluster weighting]{\textbf{Population composition across size bands, giving each cluster equal weight.} Cells show CSC-minus-background differences in the mean within-cluster category proportion, in percentage points. Each cluster contributes once, whereas Figure~\ref{fig:ch6_composition_size} weights by record counts. Blue indicates a positive difference and red a negative difference. Counts by size band are given in Table~\ref{tab:ch6_size_overview}; all displayed attributes are complete.}\label{fig:app_ch6_size_cluster_weighted}
\end{figure}

\input{assets/scotland/tab_ch6_route_sizes}
\FloatBarrier

\begin{figure}[htbp]
    \centering
    \includegraphics[width=\textwidth]{fig_app_ch6_route_composition}
    \caption[Composition by screening route]{\textbf{Composition contrasts differ between screening routes.} Cells show percentage-point differences among pooled records. Burst-only candidates (428 clusters; 18,567 records) are compared with all background (8,331 clusters; 130,021 records). Burden-only candidates (191; 2,265) and both-axis candidates (12; 1,324) are compared with burden-eligible background (3,644; 77,020). Blue indicates higher and red lower representation in the candidate group. All attributes are complete. Different comparator populations are used to respect burden eligibility; contrasts are descriptive and are not adjusted for size or epidemic context.}\label{fig:app_ch6_route_composition}
\end{figure}

\begin{figure}[htbp]
    \centering
    \includegraphics[width=\textwidth,height=0.77\textheight,keepaspectratio]{fig_app_ch6_route_entropy}
    \caption[Diversity by screening route]{\textbf{Within-cluster diversity by screening route.} Points show medians and bars the interquartile ranges, which describe variation between clusters. Left: observed entropy. Right: null-standardised entropy, with zero marked by a dotted line. Counts are shown in parentheses; all displayed entropy values are complete. Burden background is the subset of background with positive observed downstream burden. The both-axis group contains only 12 clusters and is shown descriptively. Axis scales differ by attribute on the null-standardised side.}\label{fig:app_ch6_route_entropy}
\end{figure}

\input{assets/scotland/tab_ch6_sensitivity_sizes}
\FloatBarrier

\begin{figure}[htbp]
    \centering
    \includegraphics[width=\textwidth]{fig_app_ch6_sensitivity_composition}
    \caption[Composition sensitivity to screening cutoff and size restriction]{\textbf{Descriptive composition contrasts under alternative cutoffs and size restrictions.} Cells show CSC-minus-background percentage-point differences among pooled records; blue is positive and red negative. Primary uses $p\leq0.05$ and size $\geq6$. The three alternative p cutoffs retain this size floor and redefine background at each cutoff. Size $\geq10$ and $\geq20$ retain primary labels and restrict both groups. Denominators appear in Table~\ref{tab:ch6_sensitivity_sizes}; all composition attributes are complete. Saved scores are reused without recalibration or model refitting.}\label{fig:app_ch6_sensitivity_composition}
\end{figure}

\begin{figure}[htbp]
    \centering
    \includegraphics[width=\textwidth,height=0.77\textheight,keepaspectratio]{fig_app_ch6_sensitivity_entropy}
    \caption[Diversity sensitivity to screening cutoff and size restriction]{\textbf{Entropy profiles under alternative cutoffs and size restrictions.} Points show medians and bars the interquartile ranges for CSC (blue) and background (orange) clusters. The IQR describes cluster variability, not inferential uncertainty. Definitions and group counts are given in Table~\ref{tab:ch6_sensitivity_sizes}; all displayed entropy values are complete. Observed entropy uses a common 0--1 scale; null-standardised panels have attribute-specific axes and a dotted zero line. Saved entropy values and screening scores are reused.}\label{fig:app_ch6_sensitivity_entropy}
\end{figure}

\FloatBarrier
\input{assets/scotland/tab_app_ch6_size_composition}

\clearpage
\section{Population Composition Within Clusters}\label{sec:app_ch6_descriptive}
\begin{figure}[!ht]
    \centering
    \includegraphics[width=\textwidth,height=0.72\textheight,keepaspectratio]{fig_cluster_composition_overview_demographic}
    \caption[Within-cluster composition distributions and mean contrasts]{\textbf{Within-cluster composition by CSC status.} Panels A--D show sex, age group, SIMD quintile, and urban/rural class; panel E continues on the next page. Left: proportions of member records in each category. Boxes span the interquartile range, lines mark medians, and whiskers extend to the most extreme observations within the 5th--95th percentile limits; observations beyond these limits are omitted. Right: CSC-minus-background differences in mean proportions, in percentage points (pp); blue indicates a positive difference and red a negative difference. Each of the 631 CSC and 8,331 background clusters receives equal weight, and all five attributes are complete.}\label{fig:app_composition_profiles}
\end{figure}
\clearpage
\begin{figure}[!ht]
    \ContinuedFloat
    \centering
    \includegraphics[width=\textwidth,height=0.77\textheight,keepaspectratio]{fig_cluster_composition_overview_health_board}
    \caption[]{\textbf{Within-cluster composition by CSC status (continued).} Panel E shows health board. Distributions and mean contrasts use the same definitions and equal cluster weighting as panels A--D.}
\end{figure}

\clearpage
\section{Bayesian Model Results}\label{sec:app_ch6_bayesian_results}

Figure~\ref{fig:app_random_effects_mixing} shows surveillance coefficients and policy-period and clade intercept deviations for the mixing models. The accompanying tables cover both mixing and composition models: focal attribute estimates from expanded models (Table~\ref{tab:bayesian_fixed_effects_main}), intercepts from primary and expanded models (Table~\ref{tab:bayesian_fixed_effects_intercepts}), and group-level and residual standard deviations (Table~\ref{tab:app_bayesian_random_effect_sds}).

The Chapter~6 electronic supplement, \path{chapter_06_bayesian_estimates.csv}, contains all 801 coefficient summaries for the 18 fitted models, including both specifications and entropy scales, surveillance coefficients, and individual period and clade deviations. It accompanies the thesis in the \path{supplements/chapter_06/} directory with a README defining the fields, effect scales, and diagnostic indicators. The file preserves the precision of the saved results.

\begin{figure}[htbp]
    \centering
    \includegraphics[width=\textwidth]{fig_random_effects_mixing}
    \caption[Bayesian surveillance and varying-intercept estimates for mixing models]{\textbf{Bayesian surveillance and varying-intercept estimates for cluster-level mixing models.} Points show posterior means and horizontal bars show 95\% HDIs for primary and expanded models on both entropy scales, distinguished by colour and symbol. The first three rows show surveillance coefficients per one-SD increase, included only in expanded models. Remaining rows show differences in the model baseline between policy periods and between clades. Panel A shows surveillance odds ratios and period and clade odds multipliers relative to the baseline with no group deviation; Panels B and C show differences in mean local burst and onward burden scores. Dashed lines mark one in A and zero in B and C. All models shown met the sampling criteria.}\label{fig:app_random_effects_mixing}
\end{figure}

\clearpage
% \pagestyle{plain}
\begingroup
\renewcommand{\thesisdensetablesetup}{\scriptsize\setlength{\tabcolsep}{3pt}\renewcommand{\arraystretch}{1.12}}
\input{assets/scotland/tab_bayesian_fixed_effects_main}
\endgroup
\input{assets/scotland/tab_bayesian_fixed_effects_intercepts}
\begingroup
\renewcommand{\thesisdensetablesetup}{\scriptsize\setlength{\tabcolsep}{3pt}\renewcommand{\arraystretch}{1.12}}
\input{assets/scotland/tab_app_bayesian_random_effect_sds}
\endgroup
